\section{Background and Motivation}\label{sec:background}

\subsection{Direct-to-Satellite IoT and NTN Context}
Non-terrestrial network (NTN) architectures using LEO constellations extend
coverage to remote and maritime IoT deployments
\cite{todo_ntn_iot}.  Devices communicating directly with LEO satellites
face unique上行链路 challenges: long propagation delay (order 1--10~ms),
significant Doppler shift (up to several kHz at L-band), and limited
terminal transmit power.  Location knowledge at the network side can improve
uplink scheduling, beam management, and power control, yet low-cost IoT
modules lack GNSS receivers or accurate time synchronization.

\subsection{LEO Doppler Dynamics}
A LEO satellite at altitude $h \approx 400$--1200~km moves at approximately
7.5~km/s.  The line-of-sight (LOS) range rate $\dot{\rho}(t)$ between a
ground terminal and the satellite varies from approximately
$-7.5\cos(E)$ to $+7.5\cos(E)$~km/s, where $E$ is the elevation angle.
At a carrier frequency of 1.6~GHz this produces Doppler shifts of order
$\pm 40$~kHz---well above the frequency estimation errors typical of
low-cost oscillators.

\subsection{Course Localization Utility}
Centimeter- or decimeter-level GNSS-grade accuracy is not required for many
NTN IoT use cases.  Bounding a device's position to a 100~m--10~km region of
interest (ROI) is sufficient to:
\begin{itemize}
  \item narrow beam pointing regions for satellite phased-array antennas,
  \item improve uplink timing and frequency acquisition,
  \item support anti-spoofing by detecting implausible Doppler-time signatures.
\end{itemize}
This paper targets \emph{coarse localization} within a pre-bounded ROI, not
GNSS replacement.

\subsection{Why Doppler Alone is Ambiguous}
Doppler at a single time instant depends on range-rate, which is primarily
sensitive to the radial (line-of-sight) component of terminal position.
Horizontal position components are weakly observable unless the Doppler
trajectory is tracked over an extended pass \cite{todo_doppler_localization}.
Furthermore, an unknown CFO $b_0$ and drift $b_1$ shift and tilt the
observed frequency curve, coupling oscillator state with geometry.
The \emph{shape} of the Doppler-time curve across multiple observations
provides stronger discriminative power than any single measurement, which
motivates the fingerprint approach in Section~\ref{sec:dtf}.

\subsection{Bounded ROI Assumption}
Many practical D2S IoT deployments operate in known coverage zones: a farm,
a port, an industrial site, or a known city district.  Assuming a bounded ROI
of order $10 \times 10$~km is therefore reasonable and is explicitly stated
as a precondition of the method.

\subsection{Nuisance Parameter Challenges}
Low-cost crystals exhibit CFO of tens to hundreds of Hz and drift rates of
order $10^{-2}$--$10^{-1}$~Hz/s.  These must be estimated jointly with
position; treating them as fixed biases would degrade localization accuracy
significantly.

% Placeholder citations
% TODO: fill with real references
%\cite{todo_ntn_iot}  -- 3GPP NTN IoT study
%\cite{todo_sgp4}     -- Vallado SGP4 propagator
%\cite{todo_doppler_localization} -- prior LEO Doppler localization work